一张系数表可以只是二维数组，也可以是一个线性变换在某组坐标下的完整记录——每一列写着某个基向量被搬到了哪里。选后一种读法，解方程组当场变成几何问题：目标向量 \(b\) 落不落在 \(A\) 能够到达的范围里。落在里面，解存在；不在，再怎么硬解也没有用，只能退一步去找最近的替代品。

本章沿着表示、化简、复用三步展开。

\begin{center}
\begin{tikzpicture}[font=\footnotesize,>=stealth]
  \draw[black!50,fill=blue!6] (0,2.4) rectangle (4.0,3.4);
  \node[align=center] at (2.0,2.9) {行视角\\\(m\) 张超平面交在哪里};
  \draw[black!50,fill=blue!6] (5.0,2.4) rectangle (9.0,3.4);
  \node[align=center] at (7.0,2.9) {列视角\\\(b\) 在不在列空间里};
  \draw[->,black!60] (2.0,2.35) -- (3.5,1.78);
  \draw[->,black!60] (7.0,2.35) -- (5.5,1.78);
  \draw[black!50,fill=black!4] (2.2,0.7) rectangle (6.8,1.7);
  \node[align=center] at (4.5,1.2) {消元\\换一个等价但更好读的系统};
  \draw[->,black!60] (4.5,0.65) -- (4.5,0.16);
  \draw[black!50,fill=black!4] (1.9,-0.9) rectangle (7.1,0.1);
  \node[align=center] at (4.5,-0.4) {主元落在哪些列\\唯一解 / 无穷多解 / 无解};
  \draw[->,black!60] (4.5,-0.95) -- (4.5,-1.44);
  \draw[black!50,fill=black!4] (2.4,-2.5) rectangle (6.6,-1.5);
  \node[align=center] at (4.5,-2.0) {\(PA=LU\)\\分解一次，多次求解};
\end{tikzpicture}

\emph{两种读法进来，一条化简路线出去。真正承载判断的是中间那一层：主元落在哪些列，直接决定了解是唯一、无穷多还是根本不存在；最下面一层则把这条化简路线存起来，供同一个 \(A\) 的无数个右端复用。}
\end{center}

\hyperref[sec:03-Linear-Algebra/01-Linear-Systems-and-Matrices/01]{``\texorpdfstring{从方程到矩阵：\(Ax=b\) 的三种读法}{从方程到矩阵：Ax=b 的三种读法}''} 配三副眼镜：行看约束交在哪里，列看目标能不能被现有方向拼出来，变换看矩阵把输入送去了何处。矩阵乘法的定义也在这里落地——它是变换复合在坐标下的样子，不是随手规定的搬运规则。

\hyperref[sec:03-Linear-Algebra/01-Linear-Systems-and-Matrices/02]{``消元与解的结构：从计算答案到看见约束''} 处理化简。重点不在会不会算消元，而在于消完之后主元停在哪些列：主元覆盖全部未知量则解唯一，落空一列就多一个自由度，出现 \(0=c\) 型的行则无解。全部解可以写成一个特解加上零空间，而零空间只跟 \(A\) 有关，与右端无关。

\hyperref[sec:03-Linear-Algebra/01-Linear-Systems-and-Matrices/03]{``逆矩阵与 LU 分解：撤销变换并复用消元''} 回答两个实用问题：什么样的矩阵可以撤销，以及怎样把一次消元的成本摊到许多次求解上。顺带解释为什么数值软件从不真的去算 \(A^{-1}\)，以及行交换为什么是常规动作而非应急措施。

顺序读最省力。若眼下的困惑是``为什么回归里两个特征几乎共线就会出事''，可以直接从第二节的主元与自由变量进入；若困惑是``为什么代码里从来不写求逆''，第三节的摊销账和选主元给答案。前置知识只有中学代数。

读完之后应该出现的变化是可以检验的：拿到一个线性系统，你能在不动手求解的前提下说出它的解可能是哪种形态、判断依据是什么、该盯着矩阵的哪个部位；也能说清一次分解为什么值得做，以及在什么情况下算出来的解不值得相信。本章反复出现的主元、秩、零空间、列空间四个词，此刻都还绑在消元这个具体过程上，\hyperref[ch:018]{``向量空间、基与线性变换''} 会把它们从算法里解绑。
